\title{FlashAttention-3: Fast and Accurate Attention with Asynchrony and Low-precision}

\begin{abstract}
Attention mechanisms, a foundational component of the prevalent Transformer architecture, remain a significant limiting factor in scaling large language models and facilitating applications requiring long sequences. The \textsc{FlashAttention} framework introduced optimizations to accelerate attention computations on GPUs by reducing memory access overhead. Nevertheless, this approach does not capitalize on advanced features available in the latest hardware generations; for example, \textsc{FlashAttention-2} attains merely 35\% of the theoretical throughput on H100 GPUs. In this work, we present three strategies to enhance attention performance on Hopper GPUs: utilizing asynchronous execution on Tensor Cores and TMA for (1) concurrent computation and data transfer via warp-specialization, (2) the integration of block-wise matrix multiplication with softmax evaluation, and (3) block quantization paired with incoherent execution, facilitated by FP8 low-precision hardware support. Our technique, termed \textsc{FlashAttention-3}, delivers a speedup of 1.5-2.0$\times$ on H100 GPUs, achieving up to 740 TFLOPs/s with FP16 (75\% of peak hardware efficiency), and nearly 1.2 PFLOPs/s using FP8 precision. Furthermore, we show that FP8 \textsc{FlashAttention-3} incurs 2.6$\times$ less numerical error compared to a baseline FP8 attention implementation.
\end{abstract}

\section{Introduction}
\label{sec:intro}

Within the Transformer architecture~\citep{vaswani2017attention}, the central computational challenge arises from the attention mechanism, whose calculation of self-attention scores between queries and keys results in quadratic complexity relative to sequence length. Enhancing the scalability of attention to accommodate much longer input contexts would enable advanced capabilities, such as reasoning across multiple extensive documents~\citep{guo2021longt5,shaham2022scrolls,peng2023yarn}, working with comprehensive code repositories~\citep{roziere2023code, li2023starcoder}, incorporating novel data modalities like high-resolution imagery~\citep{chen2022scaling}, audio streams~\citep{gulati2020conformer}, and videos~\citep{ho2022video}, or enabling new use cases such as tracking prolonged user interactions~\citep{sun2019bert4rec} and managing complex agent workflows extending over longer timescales~\citep{yao2022react}. These demands have spurred intensive research focused on accelerating attention computation for long sequences—using approximations~\citep{katharopoulos2020transformers,choromanski2020rethinking, tay2020efficient}, leveraging software-level optimizations~(\citep{rabe2021self,dao2022flashattention,kwon2023efficient}), and developing alternative neural architectures~\citep{peng2023rwkv,sun2023retentive,gu2023mamba}.

This paper extends the findings of \citet{dao2022flashattention}, who pioneered exact-attention algorithms by incorporating detailed knowledge of GPU hardware and execution models into their conceptual design. The work of Dao et al. \citep{dao2022flashattention} brought forward \textsc{FlashAttention}, which employs an innovative tiling method to enable attention to be processed in parallel, thus removing the need for intermediate data movement to and from slow global memory by combining all attention computations within a single GPU kernel.

In a subsequent development, \citet{dao2023flashattention2} introduced \textsc{FlashAttention-2}, reorganizing the approach to also parallelize computations across the sequence length dimension. This version executes the forward pass's internal loop over blocks of the key and value matrices, boosting GPU workload distribution and occupancy. Despite these advances, our assessment reveals that \textsc{FlashAttention-2} still falls short in resource utilization when compared to high-performance matrix-multiplication (GEMM) kernels on newer GPUs, such as achieving only 35\% utilization versus 80–90\% for GEMM kernels on the Hopper H100 GPU. This discrepancy can be partly traced to differences in low-level implementation, including the continued use of Ampere-specific instructions rather than those optimized for the Hopper architecture when deploying to Tensor Cores. Other works like ThunkerKitten~\citep{spector2024thunder} and cuDNN 9~\citep{cudnn9} have demonstrated that leveraging Hopper-optimized instructions and tile-based paradigms can accelerate attention processing while also making the software architecture more straightforward.

At a deeper level, the current design of \textsc{FlashAttention-2} relies on a straightforward synchronous execution model, deliberately omitting direct exploitation of asynchrony and low-precision arithmetic. Asynchrony arises from specialized hardware components that accelerate key computational kernels in machine learning: dedicated units such as Tensor Cores conduct matrix multiplications, while memory transfers are handled independently by modules like the Tensor Memory Accelerator (TMA), in parallel with the remainder of the CUDA cores responsible for logical and numerical tasks. The adoption of reduced-precision formats—FP8 in Hopper and FP4 in Blackwell—builds upon earlier standards such as FP16 (introduced with Pascal in 2017) and BF16 (with Ampere in 2020), and is a well-established strategy to achieve 2x or 4x improvement in throughput with no increase in power consumption or silicon footprint. Section~\cref{subsec:hardware} surveys the extended functionality made possible by Hopper in this context. 

Re-engineering \textsc{FlashAttention-2} to utilize these capabilities presents significant technical challenges: leveraging asynchrony requires orchestrating concurrent execution of matrix multiplication and softmax operations, notwithstanding the sequential dependency between them, and exploiting low-precision formats demands meticulous management of quantization errors, especially when handling outlier features, which are common in LLMs~\citep{dettmers2208llm, sun2024massive}.

For this purpose, we introduce \textsc{FlashAttention-3}, which integrates and advances three novel strategies to achieve superior performance on the latest GPU architectures.\footnote{While our empirical findings are presented based on NVIDIA's Hopper architecture, the proposed algorithm is generally applicable to any GPU platform equipped with adequate asynchronous execution and low-precision support.}

\begin{enumerate}

\item \textbf{Producer-Consumer asynchrony:} We implement a warp-specialized software pipeline that leverages asynchronous data transfers and computations on Tensor Cores by assigning producer and consumer operations to distinct warps. This design enhances the algorithm’s capacity to conceal latencies associated with memory access and instruction dispatch.
\item \textbf{Hiding softmax under asynchronous block-wise GEMMs:} We integrate low-throughput operations found in softmax—such as floating point multiply-adds and exponentials—with asynchronous WGMMA instructions for GEMM, facilitating concurrent execution. To support this, we redesign the \textsc{FlashAttention-2} approach to minimize the sequential dependency between softmax computation and GEMM steps. For instance, in the two-stage algorithmic variant, the softmax operation can process one score matrix block as WGMMA asynchronously computes the subsequent block.
\item \textbf{Hardware-accelerated low-precision GEMM:} We revise the forward computation to utilize FP8 Tensor Cores during GEMM operations, resulting in almost twice the measured TFLOPs/s throughput. To accommodate the memory layout expectations of WGMMA for FP32 accumulators and FP8 operands, we apply block quantization and incoherent processing methods. These strategies help to alleviate accuracy degradation when switching to FP8 precision.
\end{enumerate}

To empirically assess our approach, we conduct extensive benchmarks of \textsc{FlashAttention-3} on the H100 SXM5 GPU across various settings. Our results indicate: (1) the FP16 variant delivers a 1.5–2.0$\times$ acceleration in the forward phase relative to \textsc{FlashAttention-2} (with peak performance at 740 TFLOPs/s), and a 1.5–1.75$\times$ gain during backward computation; (2) FP8 achieves nearly 1.2 PFLOPs/s; and (3) for longer sequence lengths, FP16 yields superior throughput, while FP8 remains highly competitive\footnote{Specifically, for a head dimension of 64, FP8 in \textsc{FlashAttention-3} is advantageous, whereas for head dimensions of 128 and 256, it matches performance in non-causal setups and is outperformed in causal masking scenarios.} in comparison with NVIDIA cuDNN's state-of-the-art attention kernel. We further demonstrate that FP16 \textsc{FlashAttention-3} attains equivalent numerical accuracy to \textsc{FlashAttention-2} and surpasses traditional attention operators, owing to the retention of intermediate values such as softmax scaling in FP32 precision. Additionally, FP8 \textsc{FlashAttention-3}, utilizing block quantization and incoherent computation, achieves 2.6$\times$ higher accuracy than conventional attention with per-tensor quantization, particularly in settings with pronounced outlier features.

We release \textsc{FlashAttention-3} to the public under a permissive license\footnote{\textsc{FlashAttention-3}
  is available at \texttt{https://github.com/Dao-AILab/flash-attention}}, with future intentions to incorporate it into both PyTorch and Hugging Face libraries, thereby maximizing its accessibility for researchers and developers.

\section{Background: Multi-Head Attention and GPU Characteristics}
\label{sec:background}

\subsection{Multi-Head Attention}
\label{subsec:multi_head_attn}

Let $\mathbf{Q}, \mathbf{K}, \mathbf{V} \in \mathbb{R}^{N \times d}$ be the query, key and value input sequences associated to a single head, where $N$ is the sequence length and $d$ is the head dimension. Then the attention output $\mathbf{O}$ is computed as:

\begin{equation*}
  \mathbf{S} = \alpha \mathbf{Q} \mathbf{K}^\top \in \mathbb{R}^{N \times N}, \quad \mathbf{P} = \mathrm{softmax}(\mathbf{S}) \in \mathbb{R}^{N \times N}, \quad \mathbf{O} = \mathbf{P}\mathbf{V} \in \mathbb{R}^{N \times d},
\end{equation*}

The $\mathrm{softmax}$ operation is performed across each row, with the scaling parameter $\alpha$ commonly chosen as $1/\sqrt{d}$. To mitigate numerical instability that can arise with exponentiation, it is standard practice to subtract $\mathrm{rowmax}(\mathbf{S})$ from $\mathbf{S}$. In the case of multi-head attention (MHA), distinct query, key, and value projections are used for each attention head. The calculation is carried out in parallel over all heads and batches, ultimately yielding the complete output tensor.

Now let $\phi$ be a scalar loss function and let $\mathbf{d}(-) = \partial \phi / \partial (-)$ be notation for the gradient. Given the output gradient $\mathbf{dO} \in \mathbb{R}^{N \times d}$, we compute $\mathbf{dQ}$, $\mathbf{dK}$, and $\mathbf{dV}$ according to the chain rule as follows:

\begin{align*}
  \mathbf{dV} &= \mathbf{P}^\top \mathbf{dO} \in \mathbb{R}^{N \times d} \\
  \mathbf{dP} &= \mathbf{dO} \mathbf{V}^\top \in \mathbb{R}^{N \times N} \\
  \mathbf{dS} &= \mathrm{dsoftmax} (\mathbf{dP}) \in \mathbb{R}^{N \times N} \\
  \mathbf{dQ} &= \alpha \mathbf{dS} \mathbf{K} \in \mathbb{R}^{N \times d} \\
  \mathbf{dK} &= \alpha \mathbf{dS}^\top \mathbf{Q} \in \mathbb{R}^{N \times d},
\end{align*}

In this context, the relation $\mathbf{d}s = (\mathrm{diag}(p) - p p^\top)\mathbf{d}p$ holds when $p$ is defined as $\mathrm{softmax}(s)$, treating $s$ as a vector. The notation $\mathrm{dsoftmax}(\mathbf{dP})$ refers to the application of this expression to each row individually. Moreover, the backward computation for MHA can similarly be executed in parallel across both the head and batch dimensions.

\subsection{GPU hardware characteristics and execution model}
\label{subsec:hardware}

In this section, we outline the pertinent components of the GPU execution model for \textsc{FlashAttention-3}, emphasizing the specific implementation found in the NVIDIA Hopper architecture.

\paragraph{Memory hierarchy:} The organization of GPU memory follows a hierarchical structure, in which data locales offer a trade-off between capacity and bandwidth—the higher the bandwidth, the lower the capacity (\cref{tab:gpu-hierarchy})\footnote{\citet{luo2024benchmarking} measure the bandwidth of shared memory at 128 bytes per clock cycle per SM, which we scale by 132 SMs and a boost clock rate of 1830 MHz.}. At the base level, global memory (GMEM), or HBM, serves as the off-chip DRAM that all streaming multiprocessors (SMs) can access. Transfers from GMEM are automatically cached in an on-chip L2 cache. At a closer level to computation, each SM features a relatively small, programmer-controlled shared memory (SMEM) with heavy banking for efficient access. The innermost tier consists of the register file located inside each SM.

\paragraph{Thread hierarchy:} In the GPU programming paradigm, computation is structured through several layers of thread grouping. At the smallest scale are threads, which are organized into warps consisting of 32 threads each. Four consecutive warps form a warpgoup. These warpgroups are assembled into threadblocks, also referred to as cooperative thread arrays (CTAs). At a higher abstraction, threadblocks may be combined into threadblock clusters, as seen in Hopper architecture, while the largest aggregation is called a grid.

The hierarchies exhibit significant interdependence. Threads belonging to a single CTA are simultaneously executed on the same SM, while CTAs grouped within a cluster are jointly dispatched on the same GPC. SMEM is accessible to every thread in a CTA, in contrast to RMEM, which is limited to a maximum of 256 private registers per thread.

\paragraph{Asynchrony and warp-specialization:}

GPUs achieve high throughput by leveraging parallelism and asynchronous operations, which effectively mask delays caused by memory access and instruction execution. To facilitate asynchronous data transfers between GMEM and SMEM, Hopper incorporates the Tensor Memory Accelerator (TMA), a specialized hardware mechanism \cite[\S7.29]{cuda}. Additionally, Hopper’s Tensor Core—accessible through the warpgroup-level WGMMA instruction \cite[\S9.7.14]{ptx}—introduces asynchronous computation capabilities, distinguishing itself from earlier architectures like Ampere. This feature allows the Tensor Core to fetch input data directly from shared memory.

Asynchronous hardware capabilities enable warp-specialized kernels, wherein the warps within a CTA are assigned distinct producer or consumer functions, dedicating themselves exclusively to either data movement or computation tasks. This approach generally enhances the compiler’s capacity to produce efficient instruction sequences \citep{warp-specialization-2011}. Moreover, Hopper features dynamic register reallocation between warpgroups through \verb|setmaxnreg| \cite[\S9.7.17.1]{ptx}, allowing warps engaged in MMAs to utilize a greater portion of RMEM, while those handling TMA—which require only a single thread—receive less.

\paragraph{Low-precision number formats:}
\label{sec:low-precision-gpu}
Contemporary GPUs incorporate dedicated hardware modules optimized for efficient low-precision arithmetic. As an illustration, Hopper’s FP8 Tensor Cores can be exploited using the WGMMA instruction, resulting in twice the throughput per SM in comparison to using FP16 or BF16 formats.

Nevertheless, using FP8 WGMMA appropriately requires familiarity with the restrictions placed on operand layouts. Consider a GEMM operation computing $A \times B^{\top}$, where $A$ has dimensions $M\times K$ and $B$ has dimensions $N\times K$. We define an operand as \emph{mn-major} if its data is contiguous along the leading $M$ or $N$ dimension, and as \emph{k-major} if contiguity exists along the trailing $K$-dimension. In the case of FP16 WGMMA, input operands residing in SMEM can be formatted either as mn-major or k-major. Contrarily, FP8 WGMMA restricts operand layouts to k-major only. Furthermore, in workflows such as attention mechanisms—where it is desirable to chain multiple GEMMs in the same kernel—the differing memory layouts for FP32 accumulators and FP8 operands complicate the direct execution of successive FP8 WGMMA operations.

Within the scope of attention mechanisms, such layout constraints require specific adjustments to be made in the construction of an FP8 algorithm. The particular changes are outlined in \cref{sec:algofp8}.

\subsection{Standard Attention and Flash Attention} In accordance with \citet{dao2022flashattention}, we define \textbf{standard attention} as the GPU-based realization of attention that stores the intermediate matrices $\mathbf{S}$ and $\mathbf{P}$ in HBM. The central innovation of \textsc{FlashAttention} is its use of a localized softmax reduction, thereby sidestepping the need for frequent costly transfers of intermediate data and permitting the attention calculation to be conducted in a single, fused kernel. This local softmax is implemented in lines \ref{code-ws:softmax_start}-\ref{code-ws:softmax_end} of the consumer mainloop in \cref{alg:flash3_wgmma_ws_only}, in conjunction with the blockwise rescaling of $\mathbf{O}$. A straightforward derivation confirming that this approach produces the correct $\mathbf{O}$ is detailed in \cite[\S 2.3.1]{dao2023flashattention2}.

\section{FlashAttention-3: Algorithm}
\label{sec:algo}

This section presents the \textsc{FlashAttention-3} algorithm, concentrating primarily on its forward computation; the corresponding backward computation is addressed in~\cref{sec:algo_ws_bwd}. Initially, we outline the procedure for incorporating warp-specialization, supported by a circular SMEM buffer, within the foundational algorithm utilized in \textsc{FlashAttention-2}. Subsequently, we discuss how WGMMA asynchrony can be leveraged to construct a two-stage pipeline that interleaves GEMM and softmax operations. Lastly, we cover the required adjustments for FP8, detailing both layout alignment and enhancements to accuracy achieved through block quantization and incoherent processing.

\subsection{Producer-Consumer asynchrony through warp-specialization and pingpong scheduling}
\label{sec:algo_ws}

\paragraph{Warp-specialization}
Similar to the approach taken in \textsc{FlashAttention-2}, the forward computation in \textsc{FlashAttention-3} is highly parallelizable across the dimensions of batch size, number of attention heads, and query sequence length. Therefore, an overview at the CTA level sufficiently describes the algorithm, which processes a query matrix tile $\mathbf{Q}_i$ to generate its associated output tile $\mathbf{O}_i$. For clarity, the initial presentation focuses on the warp-specialization mechanism employing a circular SMEM buffer, omitting the additional aspect of GEMM-softmax overlapping. Letting $d$ denote the head dimension and $N$ represent the sequence length, we set a query block size $B_r$ that segments $\mathbf{Q}$ into $T_r = \lceil \frac{N}{B_r} \rceil$ blocks, specifically $\mathbf{Q}_1, \dots, \mathbf{Q}_{T_r}$.

\begin{algorithm}[H]
    \caption{\small\label{alg:flash3_wgmma_ws_only}\textsc{FlashAttention-3} forward pass \textbf{without} intra-consumer overlap -- CTA perspective}
    \begin{algorithmic}[1]
\REQUIRE Matrices $\mathbf{Q}_i \in \mathbb{R}^{B_r \times d}$ and $\mathbf{K}, \mathbf{V} \in \mathbb{R}^{N \times d}$ stored in HBM; key block dimension $B_c$ yielding $T_c = \lceil \frac{N}{B_c} \rceil$.
\STATE Establish pipeline to synchronize barriers with $s$-stage circular SMEM buffering.
\IF {executing in producer warpgroup}
\STATE Release reserved registers according to predetermined count.
\STATE Initiate transfer of $\mathbf{Q}_i$ from HBM to shared memory.
\STATE When load finishes, signal to consumer that $\mathbf{Q}_i$ is now available.
\FOR{$0 \le j < T_c$}
    \STATE Await confirmation that the $(j\,\%\,s)$ buffer stage has been processed by consumer.
    \STATE Start loading $\mathbf{K}_j$ and $\mathbf{V}_j$ from HBM into shared memory using the $(j\,\%\,s)$ stage.
    \STATE After completion, notify consumers that $\mathbf{K}_j$ and $\mathbf{V}_j$ are available.
\ENDFOR
\ELSE \STATE Reclaim required number of registers based on consumer warp count.
\STATE Initialize on-chip arrays: $\mathbf{O}_i = (0) \in \mathbb{R}^{B_r \times d}$; $\ell_i = (0)$ and $m_i = (-\infty)$ in $\mathbb{R}^{B_r}$.
\STATE Wait until $\mathbf{Q}_i$ is present in shared memory.
\FOR{$0 \le j < T_c$}
\STATE Wait for $\mathbf{K}_j$ to appear in shared memory.
\STATE Perform computation $\mathbf{S}_i^{(j)} = \mathbf{Q}_i \mathbf{K}_j^T$ (SS-GEMM); commit and synchronize. \label{alg:ws_only_gemm1}
\STATE Save previous $m_i$ as $m_i^{\mathrm{old}}$; update $m_i$ by taking the maximum between $m_i^{\mathrm{old}}$ and the row-wise maximum of $\mathbf{S}_i^{(j)}$. \label{code-ws:softmax_start}
\STATE Calculate $\widetilde{\mathbf{P}}_i^{(j)} = \mathrm{exp}(\mathbf{S}_i^{(j)} - m_i)$; update $\ell_i = \mathrm{exp}(m_i^{\mathrm{old}} - m_i) \ell_i + \mathrm{rowsum}(\widetilde{\mathbf{P}}_i^{(j)})$. \label{code-ws:softmax_end}
\STATE Wait for $\mathbf{V}_j$ to arrive in shared memory.
\STATE Update output: $\mathbf{O}_i = \mathrm{diag}(\mathrm{exp}(m_i^{\mathrm{old}} - m_i))^{-1} \mathbf{O}_i + \widetilde{\mathbf{P}}_i^{(j)} \mathbf{V}_j$ (RS-GEMM); commit and synchronize. \label{alg:ws_only_gemm2}
\STATE Allow producer to use the $(j\,\%\,s)$ stage again.
\ENDFOR
\STATE Finalize output by setting $\mathbf{O}_i = \mathrm{diag}(\ell_i)^{-1} \mathbf{O}_i$ and compute $L_i = m_i + \log(\ell_i)$.
\STATE Write final $\mathbf{O}_i$ and $L_i$ results to HBM as block $i$ of $\mathbf{O}$ and $L$.
\ENDIF
\end{algorithmic}
\end{algorithm}

In our deployment of \cref{alg:flash3_wgmma_ws_only} on Hopper, we employ \verb|setmaxnreg| for managing (de)allocation processes, utilize TMA to handle the loading of $\mathbf{Q}_i$ and $\{ \mathbf{K}_j, \mathbf{V}_j \}_{0 \leq j < T_c}$, and apply WGMMA within the consumer mainloop to perform the GEMMs, distinguishing the operand source via the SS or RS prefix based on whether shared memory or register file is used. To understand the operational sequence of \cref{alg:flash3_wgmma_ws_only}, it is important to realize that TMA loads operate asynchronously, permitting new loads to be issued without waiting for previous ones to complete. Additionally, throughout the producer mainloop, there are no wait operations during the initial $s$ iterations, allowing the buffer to be populated efficiently.

\paragraph{Pingpong scheduling} The combination of WGMMA and TMA's asynchronous execution and the strategy of assigning specialized tasks to warps enables the potential to concurrently execute the softmax calculation in one warpgroup while another warpgroup is carrying out GEMM. To highlight this possibility, consider that, on current hardware accelerators, matmul operations achieve significantly higher throughput than other computations, such as those involved in softmax. For instance, the H100 SXM5 GPU delivers up to 989 TFLOPS for FP16 matmul, whereas special functions like exponential computations\footnote{According to the CUDA programming guide, each streaming multiprocessor (SM) can handle 16 special function operations per clock cycle. By multiplying 16 by the total 132 SMs and the 1830 MHz clock frequency, we obtain a throughput of 3.9 TFLOPS for special functions.} required for softmax only reach 3.9 TFLOPS. During the attention forward pass in FP16, with a head dimension of 128, the demand for matmul FLOPS is 512 times greater than for exponential operations, yet the throughput for exponential is 256 times lower; consequently, the exponential computation may consume up to 50\% of the compute cycle compared to matmul. This performance disparity becomes more pronounced with FP8 precision, where matmul throughput increases while the speed for exponentials remains unchanged.

The exponential operation is handled by a distinct hardware unit—the multi-function unit—and for optimal performance, we aim to align the computation of the exponential with the period during which the Tensor Cores execute the matrix multiplication. To accomplish this, we introduce synchronization barriers (\texttt{bar.sync} instructions) that enforce an ordering: the GEMMs (GEMM1—$\mathbf{P} \mathbf{V}$ from one iteration, and GEMM0—$\mathbf{Q} \mathbf{K}^\top$ from the subsequent iteration) assigned to warpgroup 1 are scheduled ahead of those belonging to warpgroup 2. Consequently, warpgroup 1 carries out its softmax operation in parallel with warpgroup 2’s GEMM calculations. The process then alternates, with warpgroup 2 performing the softmax while warpgroup 1 returns to the GEMMs, implementing what is termed “pingpong” scheduling. This approach is depicted in~\cref{fig:pingpong_scheduling}. Although real-world scheduling with this method exhibits more complexity than the idealized version in the illustration, our experiments indicate a substantial increase in throughput (for example, improving from 570 TFLOPS up to 620–640 TFLOPS in FP16 forward passes with a head dimension of 128 and a sequence length of 8192).

\paragraph{Attention variants} In the case of multi-query attention~\citep{shazeer2019fast} and grouped query attention~\citep{ainslie2023gqa}, our implementation adopts the same methodology as \textsc{FlashAttention-2}, specifically modifying the tensor indexing so as to prevent redundant copies of $\mathbf{K}$ and $\mathbf{V}$ within HBM.

\subsection{Intra-warpgroup overlapping GEMMs and softmax}

Within a single warpgroup, it is possible to concurrently execute certain instructions from the softmax operation alongside those from the GEMMs. We outline a specific method to achieve this overlap.

In the attention algorithm, operations within the inner loop (main loop) have sequential dependencies that impede parallelization within a single iteration. For example, (local) softmax (lines \ref{code-ws:softmax_start} to \ref{code-ws:softmax_end}) relies on the output $\mathbf{S}_i^{(j)}$ of the first GEMM, while the second GEMM takes its result $\widetilde{\mathbf{P}}_i^{(j)}$ as an operand. Indeed, the wait statements in lines \ref{alg:ws_only_gemm1} and \ref{alg:ws_only_gemm2} of \cref{alg:flash3_wgmma_ws_only} serialize the execution of softmax and GEMMs. However, we can break these dependencies by pipelining across iterations through additional buffers in registers. Pursuing this idea, we propose the following two-stage\footnote{Note that the number of stages of the overlapping scheme is bounded by, but need not equal, the number $s$ of stages in the circular SMEM buffer.} GEMM-softmax pipelining algorithm:

\begin{algorithm}[H]
    \caption{\small\label{alg:flash3_wgmma}\textsc{FlashAttention-3} consumer warpgroup forward pass}
    \begin{algorithmic}[1]
      \REQUIRE  Matrices $\mathbf{Q}_i \in \mathbb{R}^{B_r \times d}$ and $\mathbf{K}, \mathbf{V} \in \mathbb{R}^{N \times d}$ stored in HBM, key block size $B_c$ with $T_c = \lceil \frac{N}{B_c} \rceil$.
      \STATE Adjust allocation of registers based on the number of consumer warps.
      \STATE Initialize, on-chip, $\mathbf{O}_i = (0) \in \mathbb{R}^{B_r \times d}$ and $\ell_i, m_i = (0), (-\infty) \in \mathbb{R}^{B_r}$.
      \STATE Await availability of $\mathbf{Q}_i$ and $\mathbf{K}_0$ within shared memory.
      \STATE Perform matrix multiplication $\mathbf{S}_{\mathrm{cur}} = \mathbf{Q}_i \mathbf{K}_0^T$ using WGMMA, commit, and synchronize.
      \STATE Free the first stage of the buffer dedicated to $\mathbf{K}$.
      \STATE Derive $m_i$, $\tilde{\mathbf{P}}_{\mathrm{cur}}$, and $\ell_i$ from $\mathbf{S}_{\mathrm{cur}}$, then apply scaling to $\mathbf{O}_i$.
      \FOR{$1 \le j < T_c - 1$}
        \STATE Ensure $\mathbf{K}_j$ is present in shared memory before proceeding.
        \label{code:mainloop_start}
        \STATE Compute $\mathbf{S}_{\mathrm{next}} = \mathbf{Q}_i \mathbf{K}_j^T$ through WGMMA; commit, no wait. \label{code:first_wgmma}
        \STATE Wait for $\mathbf{V}_{j-1}$ to become accessible in shared memory.
        \STATE Evaluate $\mathbf{O}_i = \mathbf{O}_i + \tilde{\mathbf{P}}_{\mathrm{cur}} \mathbf{V}_{j-1}$ via WGMMA; commit without waiting. \label{code:second_wgmma}
        \STATE Synchronize for WGMMA completion of $\mathbf{Q}_i \mathbf{K}_j^T$.
        \STATE Calculate $m_i$, $\tilde{\mathbf{P}}_{\mathrm{next}}$, and $\ell_i$ from $\mathbf{S}_{\mathrm{next}}$. \label{code:softmax}
        \STATE Wait for the computation $\tilde{\mathbf{P}}_{\mathrm{cur}} \mathbf{V}_{j-1}$ by WGMMA, then rescale $\mathbf{O}_i$
        \STATE Release the $(j\,\%\,s)$th and $(j-1\,\%\,s)$th stages in the buffer for $\mathbf{K}$ and $\mathbf{V}$, respectively.
        \STATE Update $\mathbf{S}_{\mathrm{cur}}$ with $\mathbf{S}_{\mathrm{next}}$.
        \label{code:mainloop_end}
      \ENDFOR \STATE Ensure $\mathbf{V}_{T_c - 1}$ is loaded in shared memory.
      \STATE Compute
      $\mathbf{O}_i = \mathbf{O}_i + \tilde{\mathbf{P}}_{\mathrm{last}} \mathbf{V}_{T_c - 1}$ utilizing WGMMA; commit and synchronize.

\STATE Finalize: Adjust $\mathbf{O}_{i}$ according to $m_{i}$. Derive $L_{i}$ using $m_{i}$ and $\ell_{i}$.
    Store $\mathbf{O}_{i}$ and $L_{i}$ in HBM, populating the $i$-th segments of $\mathbf{O}$ and $L$.
\end{algorithmic}
\end{algorithm}

\cref{alg:flash3_wgmma} serves to substitute the consumer segment in \cref{alg:flash3_wgmma_ws_only}, thereby providing the full implementation of the \textsc{FlashAttention-3} procedure for computations in FP16. Conceptually, WGMMA is employed as shorthand for the asynchronous GEMM operation. During the main processing loop, defined from lines \ref{code:mainloop_start} through \ref{code:mainloop_end}, the second WGMMA invocation in iteration $j$ (line \ref{code:second_wgmma}) is executed concurrently with the softmax calculations pertaining to iteration $j+1$ (line \ref{code:softmax}), thus enabling their overlap.

Although the pipelined architecture shown earlier has the potential to increase performance, practical implementation reveals several issues:

\paragraph{Compiler reordering} While the pseudocode assumes a sequential and ideal order of operations, the NVCC compiler may optimize by reordering instructions. Such automatic rearrangement can compromise the deliberate interleaving of WGMMA and non-WGMMA instructions intended to maximize pipeline efficiency, potentially reducing the expected speedup or causing unintended side effects. Examination of the generated SASS code confirms that the compiler does create overlapping operations, as detailed in Section~\ref{sec:2-stage-sass}.

\paragraph{Register pressure} Achieving peak throughput necessitates minimizing register spilling. However, introducing a 2-stage pipeline increases the register allocation for holding intermediate computations and maintaining state across pipeline stages. For instance, a supplementary $\mathbf{S}_{\mathrm{next}}$ must reside in registers, incurring an overhead of $B_r \times B_c \times \text{sizeof}(\text{float})$ per threadblock. The heightened register requirements could clash with efforts to enlarge block sizes—a strategy known to boost performance but also require substantial register space. Consequently, profiling is essential to assess and navigate the associated trade-offs.

\paragraph{3-stage pipelining} Building upon the 2-stage scheme, we suggest a 3-stage pipeline that aims to concurrently execute the second WGMMA and softmax operations. This enhancement could lead to greater utilization of Tensor Cores. However, this increased pipeline depth demands additional registers, intensifying the challenge of determining an optimal balance between tile dimension and pipeline stages. A comprehensive discussion of the 3-stage method and its evaluation is included in ~\cref{sec:3-stage}.

\subsection{Low-precision with FP8}
\label{sec:algofp8}

\textbf{Efficiency: layout transformations.} Executing the forward pass of \textsc{FlashAttention-3} with FP8 precision introduces unique obstacles regarding layout compatibility that are absent when using FP16.

To begin, it is important to recognize that the input tensors $\mathbf{Q}$, $\mathbf{K}$, and $\mathbf{V}$ are usually laid out so that the head dimension is the contiguous one. However, the FP8 WGMMA k-major format for the second GEMM requires that $\mathbf{V}$—more specifically, the blocks of $\mathbf{V}$ placed in SMEM—be contiguous along the sequence length dimension. Given that a TMA load cannot alter which dimension is contiguous, the solution must involve either (1) transposing $\mathbf{V}$ in global memory before computation, or (2) performing an in-kernel transpose of blocks of $\mathbf{V}$ in SMEM following their load. For option (1), this could be done by (1a) incorporating the transpose into the epilogue of an earlier stage such as rotary embedding, or by (1b) launching an independent preprocessing transpose kernel\footnote{A high-performance transpose kernel operates near the theoretical memory bandwidth of the hardware~\citep{colfax_cutlass_transpose_2024}.} to swap the strides between sequence length and head dimensions. Nonetheless, (1a) is cumbersome to adapt for general-purpose libraries, while (1b) is undesirable due to excessive memory use, particularly during inference where memory bandwidth is a constraint.

For FP8 \textsc{FlashAttention-3}, we select the second approach. Specifically, to perform the in-kernel transpose, we exploit the LDSM (\verb|ldmatrix|) and STSM (\verb|stmatrix|) instructions, which facilitate collective loading from SMEM to RMEM and storing from RMEM to SMEM by a warp, with transfers taking place at 128-byte boundaries.\footnote{The PTX reference describes LDSM/STSM as copying $8 \times 8$ matrices using 16-bit entries \cite[\S 9.7.13.4.15-16]{ptx}, but for FP8 we can pack two 8-bit values together to utilize these instructions. However, since the transpose modes of LDSM/STSM cannot separate such 8-bit packs, additional register manipulation is required between the LDSM and STSM calls to achieve a correct tile-wise transpose; the specifics are omitted here.} These LDSM/STSM operations have favorable register usage, supporting their execution within the producer warpgroup, and can alter data layouts during the copy process. Furthermore, from the second iteration onward, we can overlap the transposition of the next $\mathbf{V}$ tile with the execution of the two WGMMAs that utilize the previous $\mathbf{V}$ and the current $\mathbf{K}$ tile, effectively performing the transpose in parallel.

Secondly, we note that, contrary to FP16, the FP32 accumulator produced by FP8 WGMMA possesses a memory arrangement that does not coincide with the register-based configuration expected for operand A. \cref{fig:rmem_accum} and \cref{fig:rmem_operand} present portions of these respective layouts, illustrating how register entries are distributed per thread according to their sequential order. Employing byte permute instructions enables us to convert the first WGMMA's accumulator into a representation compatible with the next WGMMA operation and aligns with the arrangement of the $\mathbf{V}$ tile following the in-kernel transpose. Specifically, referring to \cref{fig:rmem_accum}, the sequence is reorganized as
$$\{ \verb|d0 d1 d4 d5 d2 d3 d6 d7| \},$$
with this register ordering repeated for every 8-byte segment. Regarding the $\mathbf{P}$ tile's logical structure, this transformation results in permuting the tile's columns, so that, for instance, columns $0189$ become the leading four columns. To ensure WGMMA computes the correct output tile, the in-kernel transpose can likewise be designed to generate a corresponding permutation of the rows in the $\mathbf{V}$ tile.\footnote{The flexibility provided by in-kernel transposing avoids the necessity of shuffle instructions for reassigning registers among threads, which we had previously discussed in~\citep{colfax_fp8_flashattention_2024}.}

\textbf{Accuracy: block quantization and incoherent processing.} The FP8 (e4m3) representation allocates only 3 bits for the mantissa and 4 bits to the exponent, leading to a greater accumulation of numerical error relative to formats such as FP16/BF16. Additionally, the presence of extreme outlier values in large-scale models~\citep{dettmers2208llm,
  sun2024massive}—values that significantly exceed the typical range—further complicates quantization procedures. Standard practice involves the use of per-tensor scaling~\citep{micikevicius2022fp8}, in which each tensor (for instance, $\mathbf{Q}$, $\mathbf{K}$, and $\mathbf{V}$) is associated with its own scaling factor. To mitigate the accuracy loss in FP8 attention computations, we adopt the following two strategies:

\begin{enumerate}

\item \textbf{Block quantization}: This method retains a single scalar value for each block, partitioning each tensor $\mathbf{Q}$, $\mathbf{K}$, and $\mathbf{V}$ into segments with dimensions either $B_r \times d$ or $B_c \times d$, subsequently quantizing each block independently. Notably, quantization can be combined seamlessly with certain pre-attention transformations (e.g., rotary embedding) without incurring any additional computational overhead, as rotary embedding is constrained primarily by memory bandwidth. Given that the \textsc{FlashAttention-3} algorithm operates block-wise by design, it is straightforward to rescale each segment of $\mathbf{S}$ accordingly for block quantization without extra computational cost.

\item \textbf{Incoherent processing}: To diminish the influence of outliers, we pre-multiply $\mathbf{Q}$ and $\mathbf{K}$ by a random orthogonal matrix $\mathbf{M}$ before applying FP8 quantization. The orthogonality of $\mathbf{M}$ ensures $\mathbf{M} \mathbf{M}^\top = I$, meaning the operation preserves the equivalence $(\mathbf{Q} \mathbf{M}) (\mathbf{K} \mathbf{M})^\top = \mathbf{Q} \mathbf{K}^\top$, thus leaving the computed attention values unchanged. This approach redistributes outliers because each entry in $\mathbf{Q} \mathbf{M}$ or $\mathbf{K} \mathbf{M}$ becomes a random linear combination of the corresponding original entries, effectively mitigating quantization errors. Following the strategy of \citet{chee2024quip} and~\citet{tseng2024quip}, $\mathbf{M}$ is chosen as a product of random diagonal matrices with entries $\pm 1$ and a Hadamard matrix, enabling multiplication in $O(d \log d)$ time rather than $O(d^2)$ and allowing integration with the rotary embedding step without increasing computational load.
\end{enumerate}

Empirical results confirm these approaches lower the numerical error by as much as a factor of 2.6, as discussed in \cref{sec:numerical_error}.

\section{Empirical Validation}
\label{sec:exp}

To assess the performance and precision of \textsc{FlashAttention-3}, we leverage primitives provided by CUTLASS~\citep{Thakkar_CUTLASS_2023}, including WGMMA and TMA abstractions for its implementation.

\begin{itemize}

\item \textbf{Attention performance evaluation.} We assess the execution time of \textsc{FlashAttention-3} across varying sequence lengths, conducting comparisons against the baseline PyTorch implementation, \textsc{FlashAttention-2}, the Triton-based \textsc{FlashAttention-2} variant utilizing H100-specific instructions, and a proprietary vendor’s release of \textsc{FlashAttention-2} optimized for H100 GPUs via cuDNN. Our measurements indicate that \textsc{FlashAttention-3} achieves up to 2.0$\times$ speedup over \textsc{FlashAttention-2} and is 1.5$\times$ faster compared to its Triton implementation. Furthermore, \textsc{FlashAttention-3} attains throughput up to 740 TFLOPs/s, constituting approximately 75\% of the theoretical peak on H100 hardware.
\item \textbf{Ablation experiment.} Experimental results demonstrate that enhancements in warp specialization and the GEMM-softmax pipeline substantially drive the performance gains observed in \textsc{FlashAttention-3}.
\item \textbf{FP8 attention numerical precision.} Our empirical analysis confirms that employing block quantization alongside incoherent computation reduces FP8-related numerical error in \textsc{FlashAttention-3} by a factor of 2.6$\times$.
\end{itemize}

\subsection{Benchmarking Attention}
\label{subsec:benchmark_attn}

We evaluate the execution time of various attention algorithms on an H100 80GB SXM5 GPU under multiple configurations, including the presence or absence of a causal mask and head dimensions of 64 or 128, using FP16 inputs. The findings, presented in~\cref{fig:benchmark_attn_fwd} and~\cref{fig:benchmark_attn_bwd}, reveal that \textsc{FlashAttention-3} outperforms \textsc{FlashAttention-2} with an approximate 1.5-2.0$\times$ increase in speed during the forward computation, and a 1.5-1.75$\times$ increase during the backward computation. Relative to conventional attention implementations, \textsc{FlashAttention-3} achieves speedups ranging from 3 to 16$\times$. Notably, for sequences of medium to large length (1k tokens or more), \textsc{FlashAttention-3} also exceeds the performance of a highly optimized proprietary vendor library (cuDNN), tailored specifically for H100 GPUs.

\paragraph{Benchmark settings:} We vary the sequence length as 512, 1k, ..., 16k, and set batch size so that the total number of tokens is 16k. We set the hidden dimension to 2048, and head dimension to be either 64, 128, or 256 (i.e., 32 heads, 16 heads, or 8 heads). To calculate the FLOPs of the forward pass, we use:

\begin{equation*}
  4 \cdot \text{seqlen}^2 \cdot \text{head dimension} \cdot \text{number of heads}.
\end{equation*}

Applying causal masking, we halve this quantity to reflect that roughly 50% of the elements are actually computed. For the backward pass, the FLOPs are determined by scaling the forward pass FLOPs by a factor of 2.5, since there are 2 matrix multiplications in the forward computation, whereas the backward pass requires 5 (due to the overhead from recomputation).

Additionally, we evaluate the FP8 runtime during the forward pass using comparable conditions. The performance outcomes for headdim 256 are illustrated in ~\cref{fig:benchmark_attn_fp8}, with a comprehensive set of results detailed in \cref{sec:benchmark_attn_fp8_full}.

\subsection{Ablation Study: 2-Stage Pipelining Experiments}
\label{sec:2-stage-pipelining-experiments}

We evaluate the impact of removing both the 2-stage WGMMA-softmax pipelining and warp-specialization in non-causal FP16 \textsc{FlashAttention-3}, keeping the parameters fixed at $\{\text{batch}, \text{seqlen}, \text{nheads}, \text{hdim}\} = \{ 4, 8448, 16, 128\}$. As demonstrated in~\cref{table:ablation_pipelining}, these algorithmic advancements—specifically, enabling asynchrony with warp-specialization and allowing overlap between GEMM and softmax computations—yield a substantial increase in speed, raising throughput from 570 to 661 TFLOPs.

\subsection{Numerical Error Validation}
\label{sec:numerical_error}

Given the growing focus on the numerical accuracy of \textsc{FlashAttention}~\citep{golden2024flash}, we conduct a comparison between \textsc{FlashAttention-2}, \textsc{FlashAttention-3}, and a typical attention implementation, using an FP64-based reference for evaluation. To reflect the presence of outlier features and activations as observed in LLMs~\citep{dettmers2208llm, sun2024massive}, the elements of $\mathbf{Q}, \mathbf{K}, \mathbf{V}$ are sampled from the distribution specified below:

\begin{equation*}
  \mathcal{N}(0, 1) + \mathcal{N}(0, 100) \cdot \mathrm{Bernoulli}(0.001).
\end{equation*}

Specifically, every entry is sampled from a normal distribution with a mean of zero and a standard deviation of 1; in addition, 0.1\% of the entries receive an extra independent random value drawn from a normal distribution with a standard deviation of 10. The root mean squared error (RMSE) results are presented in~\cref{table:numerical_error}. When using FP16 precision, both \textsc{FlashAttention-2} and \textsc{FlashAttention-3} reduce RMSE by a factor of 1.7 compared to the standard method, as they retain intermediate softmax computations in FP32. The FP8 baseline attention operates with per-tensor scaling, employs FP32 accumulation in matrix multiplication, and stores softmax intermediates in FP16. By leveraging block quantization and incoherent computation, \textsc{FlashAttention-3} in FP8 improves accuracy by 2.6$\times$ relative to this baseline.

\section{Dicussion, Limitations, Conclusion}
\label{sec:discussion}

Through the development of \textsc{FlashAttention-3}, we have shown that novel approaches in programming and hardware advancements—specifically asynchrony and low-precision computing—can greatly influence both the speed and fidelity of attention mechanisms. Our method achieves a 1.5-2.0$\times$ improvement in attention throughput over \textsc{FlashAttention-2}, and yields a 2.6$\times$ reduction in FP8 numerical errors compared to conventional per-tensor quantization methods. Several constraints remain, which we intend to address: enhancing the system for large language model (LLM) inference, incorporating persistent kernel architecture within the FP8 kernel,\footnote{For benchmarking purposes, the FP16 variant of \textsc{FlashAttention-3} incorporates a persistent kernel with load balancing, whereas the FP8 variant currently lacks these features. This distinction contributes to the inferior performance of FP8 \textsc{FlashAttention-3} for short sequence lengths and under causal masking, relative to FP8 cuDNN kernels.} and systematically analyzing the influence of low-precision attention during extensive-scale training. Although our evaluation centers on Hopper GPUs, we anticipate that these strategies will also be applicable to a broader set of hardware accelerators. Ultimately, we aim for robust and efficient attention mechanisms to facilitate novel possibilities, especially in tasks requiring extended contextual understanding.

\end{document}